\chapter{RFP Equilibrium}



%% RFP
% - Taylor relaxed state (need of reversal) -> magnetic shear
% - reversal can not be produced by coils ... plasma is needed
% - no need for superconducting coils and no current limit  -we don't satisfy Kruskal-Shafranov (safety factor)-
%   (10MA target to have a good condition in terms of neutrons production) -> ohmic heating is possible.
% - In RFX-mod the ohmic heating power can scale up to 60MW of coupled power (that is a huge amount of power)
%   in a stellarator to reach the same power deposition using external heating you need very high power.
% - High current possibility means -> High particle density limit ... Greenwald limit

% - why High plasma current? because 

\section{MHD, Fourier analysis, equilibrium and Modes control}
% There are several techniques to represent magnetic confinement (orbits - MHD) (brief)
% MHD is the best description to macroscopically describe the plasma as a fluid
% MHD - ideal (single double fluid)
% MHD - resistive (brief)
% Instabilities

\section{RFX-mod2 \\ \small{the opportunity of the best controllable machine}}
\cite{SONATO2003161}
\cite{doi:10.1063/1.4806765}
\cite{martin_RFX_overview}

% CHALLENGES AND SOLUTIONS IN THE DESIGN OF RFX-MOD2, A MULTI CONFIGURATION MAGNETIC CONFINEMENT EXPERIMENTAL DEVICE


% RWR
RFX-mod is a flexible \ac{RFP} toroidal device (major radius $R=2 m$ and minor radius $a=0.46 m$) with plasma current up to 2 MA \cite{2MA_RFX_Current} and volume $10 m^3$ \cite{SONATO200597}. As in all RFPs, plasma heating is purely ohmic; \acl{RFP} could in principle obtain fusion power with ohmic heating only, and with magnetic field much smaller than in a tokamak—avoiding superconducting coils. RFX-mod is equipped with a very powerful system of active coils for feedback control of plasma MHD stability: 192 coils, independently driven, cover the whole plasma surface.
% RWR
The major challenges of RFP research, and of RFX-mod in particular, are (a) rapidly advancing its performance, to assess the viability of the RFP approach to fusion; (b) providing a state-of-the-art contribution to the global task of feedback control of MHD stability, with experiments done both in RFP and in tokamak configuration; (c) focusing on the key topic of 3D magnetic shaping in a growing collaboration with the stellarator community; (d) training a new generation of fusion scientists.



\begin{figure}[ht!]
\includegraphics[width=0.5\textwidth]{img/rfx2}
\centering
\caption{RFX-mod sketch}
\label{rfx}
\end{figure}

\section{RFX diagnostics}
% kind of RFX diagnostics

%
\begin{figure}[ht!]
\includegraphics[width=1\textwidth]{img/rfx/Layout_Diagnosiche_AA10005.eps} \centering
% TODO: add figures
\caption{RFX map of diagnostics with related toroidal position.}
\label{rfx}
\end{figure}
%

\subsection{Thomson Scattering}
The Thomson scattering  has been a key-factor diagnostic at RFX since the initial operations for the study of RFP plasma configuration, yielding a fundamental correlation between the electron temperature diffusivity and the magnetic fluctuations that link the plasma confinement to the dynamo related modes~\cite{}. This also evolved in the possibility to identify a hot asymmetric island when the temperature gradient extends to the plasma core in the \ac{QSH} state.
During the first operation of RFX the diagnostic was composed by a single pulse of laser from ruby diode ( $\lambda = 694 nm$ ) passing through the plasma in the equatorial plane, and two grating photomultipliers spectrometers.
This apparatus presented some limitations: from the precision point of view it suffered of a poor quantum efficiency of the multialkali photocathode of the photomultiplier (MCP-PMT) resulting in a general lack of sensitivity and noisy profile at low electron temperatures. On the other hand the 20 chords acquisition that characterized the acquired plrofile with a 24 mm spatial resolution was unable to effectively investigate the details of the QSH magnetic islands.

Since the 2005 with the experiment upgrade to RFX-mod the diagnostic was completely renewed. In particular the replacement of the grating spectrometers with a set of four filter polycromators with avalance photodiodes (APD) provided a 30 times higher sensitivity and 84 points profile with a new narrow grained 7mm resolution ( $r/a$ from -0.96 to 0.84 ).
The new Nd:YLF laser diode ($\lambda=1053 nm$) located 15 m away from the center of the vacuum vessel, can produce an up to 7 J burst of 10 pulses per experiment shot. with a single light emission duration of 20ns $\ac{FWHM}$.

The light passing through plasma on the equatorial plane scatter in all directions and is collected by photomultipliers at three different windows at the end of short vertical ports. The 84 pairs of quarz fibers look all the plasma poloidal angles seen by magnifying optics located at the ports. All fibers can be also moved by a special mechanical support that can self orient during the experiment setup state; at the same time all optics can also automatically adjust focus to the desired radial position of interest.
A further side by side light path decomposes the acquired spectrum in 4 channels using a series of relay lenses.

The small amount of the solid angle covered by the acquisition, together with all this set of filters applied both at the input beam - aiming at reducing the stray light and maintaining the focus -, and at the output - for the selection of the plasma region of interest and for the decomposition of the energy spectrum - are the main reason for the required high power input.

This leads to a scattered set of acquisition pulses that are not continuously reconstructing the overall shape of plasma temperature but reduce ti a small set of usually 10 time events.

The recording systems is also quite complex: for RFX-mod the detected signal has been acquired by a modular 4-channels cPCI board dedicated for each of the spectrometers, acquiring data at 500 Msps each.

\subsection{Soft X Ray}

% We chose the most reliable and fast response diagnostics ... (SXR and Magnetic coils MHD)
% Other possible diagnostics can be added ( cope with time dim )
% Possibility to trigger in realtime other diagnostics ( i.e. Thomson Scattering )

\section{MHD}

\section{Newcomb SH}

\section{The complete “SCHEMA” from sensors to plasma parameters}
